(1) 设 $\theta(x,t)=T-T_a$，在半板 $0\le x\le L$ 有 $\theta_t=\alpha\theta_{xx}$，$\alpha=k/C_v$。对称中心无热流给 $\theta_x(0,t)=0$，外表面对流给 $-k\theta_x(L,t)=h\theta(L,t)$，初始 $\theta(x,0)=T_0-T_a$。

(2) 分离 $\theta=X(x)e^{-\alpha\lambda^2t}$，得到 $X''+\lambda^2X=0$。中心条件选 $X=\cos(\lambda x)$，外边界给 $k\lambda\sin(\lambda L)=h\cos(\lambda L)$，故 $\boxed{\mu\tan\mu=Bi,quad\mu=\lambda L}$，正根依次决定全部偶对称衰减频率。

(3) 最慢衰减时间 $\tau_1=L^2/(\alpha\mu_1^2)$。$Bi\ll1$ 时 $\mu_1^2\simeq Bi$，故 $\tau_1\simeq C_vL/h$，对应内部温度近似均匀的集中热容法；$Bi\gg1$ 时 $\mu_1\to\pi/2$，故 $\tau_1\to4L^2/(\pi^2\alpha)$，近似固定壁温边界。

(4) 在半板定义 $\mathcal E=\frac{C_v}{2}\int_0^L\theta^2dx$，对时间求导并分部积分得到 $\dot{\mathcal E}=-k\int_0^L\theta_x^2dx-h\theta(L,t)^2\le0$，因此非平衡温度扰动趋于消失。这个正定数学能量具有温度平方的量纲，不等于实际热量或系统热力学熵。
